% ---------------------------------------------------------------------------
% Research interests (the research/work experience timeline is commented out below)
% ---------------------------------------------------------------------------

\section{Research Interests}

\cvitem{Physics}{Dynamical horizons and strong-field black hole dynamics. Numerical relativity. Tests of general relativity. Black hole spectroscopy and ringdown. Binary black hole population inference. Third-generation and space-based detector science (Cosmic Explorer, Einstein Telescope, LISA).}
\cvitem{Methods}{Gravitational wave data analysis and Bayesian inference. Gravitational wave searches and template banks. Waveform modelling and surrogates. Pseudo-spectral methods and Cauchy-characteristic matching. High-performance computing. GPU computing. Machine learning.}
\cvitem{Publications}{13 refereed (6 first-author), 5 archived/submitted, 12 in preparation. \hyperref[sec:publications]{Full list on p.~\pageref{sec:publications}}.}

% ------------------------------------------------------------------
% Research and Work Experience: removed from the faculty CV on 2026-11-01
% (it repeated Research Interests). Kept here for grant-style CVs.
% ------------------------------------------------------------------
% \section{Research and Work Experience}
%
% \cvitem{2018 -}{\textbf{Dynamical horizons and strong-field black hole dynamics.} Quasi-local horizon geometry in binary mergers. Horizon multipoles, shear and tidal deformability. The universal structure of horizon formation. The quantitative area law and balance-law violation.}
% \cvitem{2023 -}{\textbf{Tests of general relativity and black hole spectroscopy.} GR-informed ringdown modelling. Multi-segment consistency tests. Quasi-normal mode instabilities. The apparent GR violation in GW230814. Maximum entropy in mergers.}
% \cvitem{2025 -}{\textbf{Gravitational wave data analysis and inference.} Exact accelerated time-domain likelihoods. Batched likelihood evaluation for nested sampling. Performance and convergence of stochastic samplers. Spin priors and population inference.}
% \cvitem{2023 -}{\textbf{Waveform modelling, surrogates and template banks.} Spectral numerical relativity surrogates with analytic derivatives. Massively parallel higher-order-mode template banks for precessing binaries. Eccentric and precessing binary simulations.}
% \cvitem{2018 -}{\textbf{Numerical relativity methods and infrastructure.} Cauchy-characteristic matching on compactified null slices. Pseudo-spectral methods and numerical PDEs. Apparent horizon finding in SpEC and SpECTRE. Compile-time optimization and HPC benchmarking. Comparison infrastructure for NR catalogues.}
% \cvitem{2025 -}{\textbf{Next-generation detector science.} Science traceability and tests of general relativity for Cosmic Explorer. Forecasting for third-generation and space-based detectors.}
% \cvitem{2018 - 2021}{\textbf{Scientific computing and visualization.} Parallel computing and data I/O on HPC systems. 3D scientific visualization of numerical relativity data.}
% \cvitem{2015 - 2020}{\textbf{Earlier work.} Causal structure of spacetime, cosmic censorship and chronology protection. Black holes in higher dimensions. Scattering of ringdown signals from black holes surrounded by matter. Magnetohydrodynamics and galactic dynamos. Cosmological parameter estimation from supernovae.}

% ------------------------------------------------------------------
% Condensed from the previous 33-item chronological list, preserved below.
% ------------------------------------------------------------------
% \cvitem{2026-}{GR-informed ringdown modelling}
% \cvitem{2026}{Cauchy-characteristic matching via generalized-harmonic evolution on compactified null slices}
% \cvitem{2026}{Numerical relativity surrogate waveform modelling with action-angle parameterisation (NRHJSur3dq8)}
% \cvitem{2026}{Massively parallel higher-order-mode template bank generation for precessing binaries}
% \cvitem{2026}{Universal structure of horizon formation in generic binary black hole mergers}
% \cvitem{2026}{Balance-law violation of finite-length waveforms and the geometry of precession}
% \cvitem{2026}{Multi-segment consistency tests of general relativity}
% \cvitem{2026}{Maximum entropy in binary black hole mergers}
% \cvitem{2025}{Spin priors for GW astronomy}
% \cvitem{2025}{Investigation of apparent violation of GR with GW230814}
% \cvitem{2025}{Area Law analysis for GW250114}
% \cvitem{2025}{Studies of performance and convergence of stochastic sampling in Gravitational Wave parameter estimation}
% \cvitem{2025}{Accelerated time domain analysis in gravitational wave astronomy}
% \cvitem{2025}{Science Traceability and Tests of General Relativity with Cosmic Explorer}
% \cvitem{2025}{Comparison infrastructure for Numerical Relativity Catalogs}
% \cvitem{2023}{Numerical Simulations of Eccentric, Aligned/Precessing BBH systems}
% \cvitem{2023}{An accelerated Apparent Horizon Finder based on dynamical horizon geometric structure}
% \cvitem{2023}{Investigation of QNM instabilities using different numerical methods}
% \cvitem{2023}{Horizon finding with SpECTRE code}
% \cvitem{2022}{Dynamical Horizons with SpEC}
% \cvitem{2022}{Pseudo Spectral methods and numerical PDE solving}
% \cvitem{2022}{Compile time optimizations of SpEC NR codebase ICTS-TIFR}
% \cvitem{2022}{Benchmarking HPC at ICTS-TIFR, SpEC}
% \cvitem{2021}{3D scientific visualizations}
% \cvitem{2020}{Numerical solutions to scattering of ringdown signals from black holes surrounded by matter}
% \cvitem{2019}{Cosmic Censorship and Chronology Protection problems}
% \cvitem{2018}{Numerical Relativity using The Einstein Toolkit}
% \cvitem{2018}{Physics of Dynamical Horizons using Numerical Relativity}
% \cvitem{2018}{Black holes in Higher dimensions}
% \cvitem{2018}{Data analysis, including Gravitational wave data analysis}
% \cvitem{2017}{Magneto Hydro Dynamics and the Turbulent dynamos with application to galactic dynamos}
% \cvitem{2015}{Estimation of Cosmological Parameters using Super-Novae data}
% \cvitem{2015}{Causal structure of spacetime and Chronology Protection}
